\uuid{wBs9}
\exo7id{7322}
\titre{exo7 7322}
\auteur{mourougane}
\organisation{exo7}
\datecreate{2021-08-10}
\isIndication{false}
\isCorrection{false}
\chapitre{Groupe, anneau, corps}
\sousChapitre{Anneau}
\cpge{mpsi-mp2i.polynomes-fractions-rationnelles.c, pcsi.polynomes.c, ptsi.polynomes.c, mpsi-mp2i.polynomes-fractions-rationnelles.b, pcsi.polynomes.b, ptsi.polynomes.b, mpsi-mp2i.nombres-reels-suites.a, pcsi.raisonnement-vocabulaire-ensembliste.c, ptsi.raisonnement-vocabulaire-ensembliste.c, bcpst1.polynomes-reels.b, bcpst1.polynomes-reels.a, bcpst2.polynomes.b}
\module{Algèbre}
\niveau{L2}
\difficulte{}

\contenu{
\texte{
Soit $P = a_0 + a_1X + \cdots + a_nX^n \in \Z[X]$. On veut déterminer si $P$ a des racines rationnelles.
}
\begin{enumerate}
    \item \question{On suppose que $P$ a une racine rationnelle non nulle $x$, avec $x = \frac{p}{q}$ et $\pgcd(p,q) = 1$. Montrer que $p$ divise $a_0$ et $q$ divise $a_n$.}
    \item \question{Le polynôme $7X^3 - 5X^2 -9X +4$ a-t-il des racines rationnelles ? et $X^4 - 2X^2 - 3$ ?}
    \item \question{Soit $n \in \N^*$. Montrer que $\sqrt{n}$ est soit un entier, soit un irrationnel.}
\end{enumerate}
}
